\documentclass[twoside,11pt]{article}

% Set some margins (columnsep is used just for 2-col layouts) and
% the paper size.
\usepackage[letterpaper, margin=0.75in, columnsep=0.5in]{geometry}

% Set up fonts. Use the Bembo-for-poor-people typeface for text, Libertine
% for Math equations and Inconsolata for monospace (you cannot use anything
% else with pdfLatex). Note the number styling (`osf`).
\usepackage[T1]{fontenc}
\usepackage[expansion=false]{microtype}
\usepackage[p]{fbb}
\usepackage{libertinust1math}
\usepackage{amsmath, amsfonts, amsthm}
\usepackage{inconsolata}

% Other packages
\usepackage[hang,flushmargin]{footmisc}
\usepackage[labelfont=bf, textfont=it]{caption}
\usepackage{adjustbox}
\usepackage{array}
\usepackage{arydshln}
\usepackage{blindtext}
\usepackage{caption}
\usepackage{color,soul}
\usepackage{enumitem}
\usepackage{float}
\usepackage{framed}
\usepackage{graphicx}
\usepackage{hyperref}
\usepackage{minted}
\usepackage{multicol}
\usepackage{needspace}
\usepackage{pdflscape}
\usepackage{rotating}
\usepackage{subcaption}
\usepackage{tikz}
\usepackage{titlesec}
\usepackage{titling}
\usepackage{xcolor}
\usepackage{xfp}

% References. Managed by BibLaTeX with Biber as the backend for APA, MLA, etc.
% Hyperref will make the citations clickable.
% https://www.overleaf.com/learn/latex/Biblatex_bibliography_styles
\usepackage[style=numeric,backend=biber,sorting=none]{biblatex}

% =============================================================

% Customize maketitle
\pretitle{\begin{center}\Huge} \posttitle{\end{center}}
\preauthor{\begin{center}\Large\rmfamily} \postauthor{\end{center}}
\predate{\begin{center}\small} \postdate{\end{center}}

% % Patch Tabular to use monospaced font for the tables
% \AtBeginEnvironment{tabular}{\sffamily}

% Add circles to things. Good for inline lists. \circled{1}, for example.
\newcommand*\circled[1]{\tikz[baseline=(char.base)]{\node[shape=circle,draw,inner sep=1pt] (char) {#1};}}

% Caption setup. Make the labels smallcaps and bold. Make the entire font size
% smaller. Rename the figure and table names.
\captionsetup{
  labelfont={bf,smaller,sc},
  textfont=smaller,
  labelsep=period,
  figurename=FIGURE,
  tablename=TABLE
}

% Set some margins, set the typeface, do not indent paragraphs,
% adjust paragraph line heights.
\setlength{\parindent}{0em}
\setlength{\parskip}{0.75em}

% Customize footnotes
\setlength{\footnotesep}{\baselineskip}
\setlength{\footnotemargin}{0.75em}
\setlength{\skip\footins}{1em}

% Customize lists
\setlist[itemize]{itemsep=0.5pt, topsep=0pt, leftmargin=12pt}
\setlist[enumerate]{itemsep=0.5pt, topsep=0pt, leftmargin=12pt}

% Independent sign
\newcommand\independent{\protect\mathpalette{\protect\independenT}{\perp}}
\def\independenT#1#2{\mathrel{\rlap{$#1#2$}\mkern2mu{#1#2}}}

% Customize title separation
% https://tex.stackexchange.com/a/108747
\titlespacing*{\section}{0pt}{1pt}{1pt}
\titlespacing*{\subsection}{0pt}{1pt}{1pt}
\titlespacing*{\subsubsection}{0pt}{1pt}{1pt}

% Header and footer
\usepackage{fancyhdr}
\pagestyle{fancy}
\fancyhf{}
\fancyhead[L]{\color{lightgray}\fontsize{8pt}{1em}\selectfont\uppercase{Research Methods}}
\fancyhead[C]{\color{lightgray}\fontsize{8pt}{1em}\selectfont{Spring 2026}}
\fancyhead[R]{\color{lightgray}\fontsize{8pt}{1em}\selectfont\thepage}
% \fancyfoot[R]{\color{lightgray}\fontsize{8pt}{1em}\selectfont\thepage}
\renewcommand{\headrulewidth}{0pt}

% Add bibliography
\addbibresource{References.bib}

\definecolor{phendopink}{HTML}{B73275}

% =============================================================

\begin{document}
\raggedbottom

\title{Bridging Conversational Expression \textit{and} Computational Health Representation \textit{in} Endometriosis}
\author{
  Nikhil Anand
}
\date{
  {\scriptsize RESEARCH METHODS, SPRING 2026}\\[0.5em]
  Submitted \today
}

\maketitle

\noindent\rule{\columnwidth}{0.25pt}
\vspace{0.0625em}

\begin{tikzpicture}[remember picture,overlay,shift={(current page.north east)}]
  \node[anchor=north east,xshift=0,yshift=0]{\includegraphics[width=7.2em]{./media/qr-code.png}};
\end{tikzpicture}

% =============================================================

\section*{\centering Abstract}

Endometriosis is a chronic, inflammatory, enigmatic condition affecting 6--10\% of women of reproductive age. Patients carry complex, embodied, often inexpressed experience of their disease, and communicate with their providers across a temporal-resolution gap: patients live their condition moment-to-moment while providers operate visit-to-visit. Existing self-tracking apps require patients to distill that experience into structured forms, with documented mistranslation. Existing timeline visualizations are designed for clinicians and analysts. Existing computational phenotypes have been shown to misalign with patients' own self-assessments.

This report describes \textit{EndoChron}, a mobile-first prototype that bridges the gap between conversational expression and computational health representation in endometriosis. Patients record unstructured, spoken accounts of their moment. On-device transcription and extraction map their speech onto the Phendo citizen-science taxonomy. An editable, multi-resolution timeline materializes from the accumulated speech, in addition to a body-map view that produces a patient-curated, visit-preparation artifact for a clinical encounter.

Design is grounded in formative interviews with five participants (three with endometriosis), in qualitative literature on endometriosis communication, and theoretically in Distributed Cognition and Coiera's Communication-Conversation model. A working prototype is delivered and may be accessed at \href{https://docker.nikhil.io}{\texttt{docker.nikhil.io}} with the source available on \href{https://github.com/afreeorange/endochron}{GitHub}. A Phase~II RCT is proposed to test clinical efficacy of the completed solution, with PEPPI-5 (Perceived Efficacy in Patient--Physician Interactions) as the primary outcome, in a stratified two-arm design that recruits from the existing Phendo cohort.

\vfill

\begin{figure}[H]
  \centering
  \makebox[\textwidth][c]{%
    \includegraphics[width=0.4\textwidth]{./media/logo.jpg}
  }
\end{figure}

\newpage

% =============================================================

\renewcommand*\contentsname{\centering Table \textit{of} Contents}
{\footnotesize
  \tableofcontents
}
\newpage

% =============================================================

\section{Introduction \& Significance}

Endometriosis is a chronic and inflammatory condition characterized by the growth of endometrial-like tissue outside the uterus\cite{Zondervan2018Endometriosis} that affects approximately 6--10\% of women of reproductive age worldwide\cite{Zondervan2018Endometriosis}. It is enigmatic in its etiopathogenesis, with no known biomarker or cure, no standards for treatment, and heterogenous symptomatic presentation\cite{Takeuchi2024Endometriosis,McKillop2019Thesis}. The diagnostic delay for endometriosis averages 7--11 years from symptom onset\cite{Agarwal2019Endometriosis,Soliman2017Endometriosis,staal2016diagnostic}. On average, patients see 4-6 clinicians before receiving a confirmation of their condition\cite{McKillop2019Thesis}.

\subsection{Communication and Management}

We argue that the personal \textit{and} clinical experience of endometriosis are shaped by a single underlying difficulty in two related ways: patients struggle to make sense of an embodied condition on their own, and struggle to convey this experience to providers in the circumscribed time permitted by a clinical encounter. This is evidenced by qualitative research that has consistently documented disconnects and misalignments in how patients and providers understand, represent, and discuss the condition\cite{Pichon2021Divided,Armour2019Endometriosis,Mardon2023EndometriosisReview,OHara2019EndometriosisReview,Denny2009Endometriosis,Young2015Endometriosis}. In light of these difficulties, we identify three allied concerns that inform the goals of this project:

\begin{enumerate}
  \item \textbf{Patients live with day-to-day uncertainty and carry complex, embodied, often inexpressed or underexpressed experience.} The day-to-day experience of endometriosis includes pain that varies in location, character, and intensity, gastrointestinal and urinary symptoms, fatigue, mood disturbance, sexual dysfunction, and functional limitation across several dimensions of health and activities of daily living\cite{McKillop2019Thesis,Urteaga2020Phenotypes}. Much of this disruptive experience is filtered through stigma, social isolation, and a history of dismissal by caregivers, which leads patients to under-report or minimize what they experience\cite{Ballard2006Endometriosis,Pichon2025Betrayal}\footnote{``\textit{It becomes its own full-time job.}''\cite{Pichon2021Divided}}.

  \item \textbf{Patients and providers operate at different temporal resolutions.} Patients live day-to-day or moment-to-moment with their condition, whereas providers think visit-to-visit. The result is that patients must compress weeks or months of embodied experience into the short and circumscribed length of of a clinical visit, whilst simultaneously recognizing and managing the emotional and cognitive labor of self-advocacy\cite{Pichon2021Divided}\footnote{``\textit{It feels like a trial, I have seven minutes as a lawyer to prove my case, to present enough evidence, to hit the particular buzzwords.}''\cite{Pichon2021Divided}.}.

  \item \textbf{Existing computational representations of the patient experience misalign with the patient's own self-assessment.} Work on Phendo, a specialized self-tracking app for endometriosis, has shown that even rigorously learned computational phenotypes match patients' own assessments of their health state only 18\% of the time, and rule-based phenotypes match only 35\% of the time\cite{Pichon2025Betrayal}. Data representations built from their patients' own self-reports fail to align with narratives of patients' embodied experiences.
\end{enumerate}

\subsection{Current Informatics Interventions}

Three classes of informatics interventions address many dimensions of this overall tension but each leaves a critical and unaddressed gap.

\begin{enumerate}
  \item \textbf{Self-tracking apps} like Phendo\cite{McKillop2019Thesis,CitizenEndoPhendo}, Bearable\cite{BearableApp}, and Visible\cite{VisibleHealth}\footnote{Phendo being an endometriosis-specific app. The others are mentioned as `general-purpose', `popular'  apps, popularity being determined by cursory examinations of both Apple App Store downloads \& searches across \texttt{MyEndoMetriosisTeam.com}} enable rich symptom data collection across many disease dimensions. They have engaged patients longitudinally\cite{Urteaga2022Engagement} supported computational phenotyping\cite{Urteaga2020Phenotypes,McKillop2019Thesis}. However, they require patients to distill the complexity of their embodied, emotional experience into structured fields (typically multi-part forms that primarily entail checkboxes, radio and other buttons, and text inputs) which yield mistranslations between the complexity and `messiness' of embodied patient experience and what an informatics intervention may utilize for downstream analysis apropos sensemaking, insight generation, or visit preparation.
  \item This under/misrepresentation suffuses \textbf{Timeline Visualizations} which do support analyst and provider decision-making processes but do not incorporate the lived experience of the patient and/or their narrative\cite{Ledesma2019HealthTimeline}.
  \item \textbf{Conversational AI for disease self-management} is emergent\cite{kocielnik2018reflection} but has not been applied to or studied in chronic or enigmatic conditions, nor has it been integrated with temporal or anatomical visualizations that support clinical communication.
\end{enumerate}

\subsection{The Gap}

We posit that \textit{there is currently no informatics intervention that maps a patient's spoken, `messy', embodied account of an enigmatic and chronic condition like endometriosis to rich and expressive self-authored, multi-temporal artifacts for reflection, sensemaking, and shared clinical decision-making.}

% =============================================================

\section{Innovation}

We suggest an intervention, \textit{EndoChron}, that addresses this gap by employing \textbf{unstructured voice as the primary input modality} of richer and authentic expression of lived experience\footnote{With the express understanding that transcription loses the import of, for example, sarcasm and/or volume.} to bridge health expression between the conversational and the computational in enigmatic and chronic conditions that emerges \textit{patient authored} and \textit{clinically useful} artifacts.

Three secondary innovations follow: \circled{1} the use of an LLM to map free-form patient speech onto an established, citizen-science-derived clinical taxonomy (Phendo) rather than an investigator-defined ontology, \circled{2} the materialization of that taxonomy-mapped speech as an \textbf{editable multi-resolution timeline}, offered with initial, temporally contextual and condition-based LLM-based summaries that the patient entirely controls, and \circled{3} the mapping of longitudinal speech-derived data onto an anatomically organized visit-preparation artifact, \textbf{a body map}, that the patient explicitly curates for disclosure to a clinician.

\subsection{Comparison with Prior Tools}

Table~\ref{tab:prior-tools} situates EndoChron against the four most relevant classes of prior work, on four dimensions: input modality, the author of the representation, the temporal resolution supported, and whether the tool produces an artifact designed for the clinical visit.

\begin{table}[htbp]
  \centering
  \scriptsize
  \renewcommand{\arraystretch}{1.5}

  \makebox[\textwidth][c]{%
    \begin{tabular}{
        >{\raggedright\arraybackslash}p{3.5cm}
        >{\raggedright\arraybackslash}p{3.0cm}
        >{\raggedright\arraybackslash}p{2.7cm}
        >{\raggedright\arraybackslash}p{3.0cm}
        >{\raggedright\arraybackslash}p{3.5cm}
      }
      \hline
      \textbf{Tool class} & \textbf{Input modality}                        & \textbf{Authorship} & \textbf{Temporal resolution} & \textbf{Visit-prep artifact} \\
      \hline

      Phendo, Bearable, Visible, other self-tracking apps
                          & Structured forms
                          & System/Designer-authored from patient inputs
                          & Primarily Day-level
                          & None patient-curated                                                                                                               \\
      \hline

      Clinical timeline visualizations \cite{Ledesma2019HealthTimeline,Rind2013InteractiveEHR}
                          & Structured EHR data
                          & Clinician/analyst-authored
                          & Visit-to-visit
                          & For clinician, not patient                                                                                                         \\
      \hline

      Computational phenotyping from PGHD \cite{Urteaga2020Phenotypes,McKillop2019Thesis}
                          & Self-tracked data from app
                          & Algorithm-authored
                          & Aggregated over study period
                          & None (research output)                                                                                                             \\
      \hline

      Conversational AI for chronic disease
                          & Voice / text
                          & Mixed, AI summary of session
                          & Session-level
                          & Not visit-curated                                                                                                                  \\
      \hline

      \textbf{EndoChron (this work)}
                          & \textbf{Unstructured Speech}
                          & \textbf{Patient-authored, AI-assisted}
                          & \textbf{Moment $\rightarrow$ Year}
                          & \textbf{Body map + Narrative, Patient-Curated}                                                                                     \\
      \hline
    \end{tabular}%
  }

  \caption{EndoChron in relation to prior tool classes. ``Patient-authored'' means that the patient, not the system, is the final authority on what the representation says\cite{Pichon2025Betrayal}.}
  \label{tab:prior-tools}
\end{table}

The novelty of EndoChron is not any single column of this table but rather a combination thereof. No prior tool we are aware of accepts unstructured speech, preserves patient authorship through every downstream representation, supports temporal resolutions from moment to year, and produces a visit-preparation artifact that the patient explicitly controls.

\subsection{A Tractable Problem}

Three recent, converging developments make this intervention possible at the time of this writing. \circled{1} On-device speech transcription via small LLMs (Whisper-class models) has reached production quality on consumer smartphones, enabling the privacy guarantee that patient audio never leaves the device\cite{Radford2023Whisper}. \circled{2} The Phendo program has produced a validated citizen-science taxonomy of endometriosis signs and symptoms\cite{McKillop2019Thesis} that can serve as a controlled vocabulary for the extraction step, replacing the need for investigator-defined ontologies. \circled{3} Qualitative research on patient--provider misalignments in endometriosis has matured to the point where the design requirements for a patient-authored intervention are well understood\cite{Pichon2021Divided,Pichon2025Betrayal}.

% =============================================================

\section{Approach}

\subsection{Preliminary Work and Requirements Gathering}

To ground EndoChron's design in lived experience, we conducted formative sessions remotely with five participants assigned female sex at birth, three with a confirmed endometriosis diagnosis. Table \ref{tab:participant_demographics} lists some demographic characteristics. 1-2 sessions were conducted per participant at this phase, each lasting an average of $\approx$50 minutes.

Sessions were semi-structured and used a master-apprentice framing\cite{Beyer1997ContextualDesign}. They followed contextual inquiry conventions where participants walked us through their current practices for \circled{1} noticing and recording symptoms day-to-day (three participants used at least one app, one journaled, and one did not engage in any self-tracking), \circled{2} making sense of patterns over time, and \circled{3} preparing for clinical encounters.

\begin{table}[h]
  \centering
  \footnotesize
  \begin{tabular}{|c||c|c|c|c|l|}
    \hline
    \textbf{Participant} & \textbf{Age} & \textbf{Endometriosis Dx} & \textbf{Age at Dx} & \textbf{Uses Self-Tracking App(s)} & \textbf{Other Conditions} \\
    \hline
    \hline
    A                    & 53           &                           & --                 & \checkmark                         & Anxiety, Depression       \\
    \hline
    B                    & 34           & \checkmark                & 26                 & \checkmark                         & Pre-diabetic              \\
    \hline
    C                    & 38           &                           & --                 &                                    &                           \\
    \hline
    D                    & 46           & \checkmark                & ``30s''            &                                    &                           \\
    \hline
    E                    & 39           & \checkmark                & 31                 & \checkmark                         & Hashimoto's, Depression   \\
    \hline
  \end{tabular}
  \caption{Participant demographics, self-tracking app usage, and other conditions during Formative Design.}
  \label{tab:participant_demographics}
\end{table}

All sessions were recorded and transcribed with participant consent. All personally identifying information was deleted from all records immediately after analysis\footnote{Our Human Subjects Research statement may be found in Section~\ref{sec:human-subjects}.}. Thematic Analysis of the transcripts yielded the following, broad, recurring patterns:

\begin{enumerate}
  \item \textbf{Input Friction.} Participants described the act of opening a self-tracking app, navigating to the right view/field, and selecting an option as itself an ``\textit{annoying}'' barrier to journaling. One of the three endometriosis participants reported having abandoned self-tracking apps for this reason\footnote{Primarily. Privacy was an equal concern in her case.}. One of two non-endometriosis participants evinced lower annoyance with this input modality.
  \item \textbf{Mists of Time.} All participants described difficulties with recalling not just symptoms but the surrounding context (food, sleep, life events, for instance) when preparing for a clinical visit. Two of three participants who used at least one self-tracking app reported that the visualizations offered were ``\textit{okay}'' and ``\textit{kinda useful}''. The sole endometriosis participant who used a self-tracking app (not specific to her condition) reported difficulty with trying to connect her app's graphical summaries to her flare-ups.
  \item \textbf{Pre-visit Cognitive and Emotional Load.} Four of five participants, including all three with endometriosis, described a pattern of preparing for clinical visits the day before, trying to construct a coherent narrative (``\textit{hit the key points}'') from memories and/or self-tracking data. ``\textit{You're just really worried about what \underline{you're} going to say and what \underline{they're} going to say.}''
  \item \textbf{Loss of Authorship.} Two of three endometriosis participants described instances where a provider summarized their experience to them in a way that felt ``\textit{incongruent}.'' One of them reported correcting their provider at earlier clinical visits (``\textit{He gets me now.}'')
  \item \textbf{Privacy Concerns.} Four of five participants reported concerns with how their data was gathered and used.
\end{enumerate}

Our findings are consistent with the broader qualitative literature on endometriosis communication\cite{Pichon2021Divided,Pichon2025Betrayal,Mardon2023EndometriosisReview} and directly motivate the design goals enumerated in Section~\ref{sec:design-goals}.

\subsection{Flow Model of Endometriosis Communication}

Following Holtzblatt and Beyer's Contextual Design methodology\cite{Beyer1997ContextualDesign}, we modeled the work of communicating about endometriosis as a Flow Model with three principal \textit{entities}, the \textbf{patient}, the \textbf{device}, and the \textbf{provider}, and three \textit{activities} through which health representations propagate: \textbf{capture}, \textbf{reflection / sensemaking}, and \textbf{clinical communication}. While not a Flow Diagram, the Architecture in Figure \ref{fig:architecture} labels all entities and activities.

With current self-tracking applications, and using a simple Flow Model, several gaps may be noted:

\begin{itemize}
  \item Between Embodied Experience $\rightarrow$ Capture, structured-input modalities introduce under/mistranslation.
  \item Between Capture $\rightarrow$ Reflection, existing apps offer limited sensemaking.
  \item Between Reflection $\rightarrow$ Clinical Communication, no artifacts exist in a form the patient owns and curates.
\end{itemize}

\subsection{Theoretical Grounding}

The design of EndoChron is grounded in Distributed Cognition, which maintains that cognitive tasks are not accomplished by an individual mind in isolation but by a socio-technical system in which representations propagate through representational media\cite{furniss2019sociotechnical}.

The clinical visit for an endometriosis patient fits with this theory: months of embodied experience must be `moved' from the patient's body and memory, through some representational medium, into a form the provider can perceive and act on. The gaps we observed in formative interviews occured at the interfaces between the body $\rightarrow$ form, form $\rightarrow$ clinician, clinician's summary $\rightarrow$ patient's confirmation.

EndoChron is designed as a chain of representational media (speech $\rightarrow$ transcript $\rightarrow$ multi-resolution timeline $\rightarrow$ body-map summary) intended to reduce the cognitive load on the patient at the moment of the visit \textit{and} to preserve patient authorship at each transformation. It aims to externalize clinically relevant experiences into persistent visual and linguistic representations, reducing reliance on memory-intensive tasks.

A secondary theoretical grounding is Coiera's Communication-Conversation Model\cite{Coiera2000ConversationModel}, where the patient and provider need to establish \textit{common ground} for clinical work to proceed to the satisfaction of both patient and provider. EndoChron's pre-visit artifact, the Body Map, is intended to be solid ground (established before the visit) that permits limited visit time to be allocated to the \textit{shifting-ground} work for shared decision making.

\subsection{Design Goals}
\label{sec:design-goals}

Four design goals follow from the formative interviews and theoretical commitments described:

\begin{enumerate}[label=\textbf{DG\arabic*.},leftmargin=*,labelsep=0.5em]
  \item \textbf{Eliminate input friction by accepting unstructured speech as the primary input modality.} The patient's account of their experience should be the input to the system, not a translation of that account into checkboxes. Audio is transcribed on-device and discarded. The patient never sees a structured form during capture.
  \item \textbf{Preserve patient authorship across every downstream representation.} The patient must be allowed to inspect, correct, or reject what the system has inferred at every step in the transformation. The system must be positioned as a ``sounding board,'' not an authority\cite{Pichon2025Betrayal}
  \item \textbf{Support sensemaking at the temporal resolution that matches lived experience.} The system must support navigation from the moment level through days, weeks, months, and years. This addresses both the day-to-day resolution at which patients live their condition and the visit-to-visit resolution at which thier providers operate\cite{Pichon2021Divided}.
  \item \textbf{Produce a visit-preparation artifact that the patient controls.} The system must transform longitudinal speech-derived data into a clinically usable representation, that may span arbitrary time periods, that the patient can curate, edit, and annotate with what to disclose at their visit/encounter.
\end{enumerate}

Table~\ref{tab:needs-goals-features} maps each goal to the user need it addresses and to the EndoChron features that implement it. The full RFC-2119 requirements spec may be found in Appendix~\ref{sec:requirements}.

% Add to preamble, after \usepackage{xcolor}
\definecolor{phendopink}{HTML}{B73275}  % darker rose, matches Phendo brand

\begin{table}[H]
  \centering
  \footnotesize
  \renewcommand{\arraystretch}{1.5}
  \begin{tabular}{p{0.33\textwidth}|p{0.3\textwidth}:p{0.33\textwidth}}
    \hline
    \textbf{User Need / Gap}                                                                                                                                             & \textbf{Design Goal}                                                                        & \textbf{EndoChron Features}                                                                                                                                                                                                                                                                                                                      \\
    \hline
    \hline
    Patients abandon self-tracking because structured forms create friction and mistranslate embodied experience (Interviews; \cite{McKillop2019Thesis}).                & \textbf{DG1.} Accept unstructured speech as the primary input modality.                     & \textbf{Recording UI:} Voice recording interface, on-device transcription, no structured input imposed \textcolor{phendopink}{(Capture~1--3)}.                                                                                                                                                                                                   \\
    \hline
    Patients distrust algorithmic representations of their experience after experiences of dismissal and misalignment\cite{Pichon2025Betrayal,Ballard2006Endometriosis}. & \textbf{DG2.} Preserve patient authorship across every representation.                      & Transparent display of extracted Phendo terms \textcolor{phendopink}{(Capture~6)}, patient override of all inferences including day rating \textcolor{phendopink}{(Capture~7)}, inline editing of summaries \textcolor{phendopink}{(Reflective Sensemaking~6)}, patient-curated disclosure in visit prep \textcolor{phendopink}{(Visit Prep~6)}. \\
    \hline
    Patients cannot recall context surrounding flare-ups in the visit \& need sensemaking support at multiple temporal scales (Interviews; \cite{Pichon2021Divided}).    & \textbf{DG3.} Support sensemaking at the temporal resolution that matches lived experience. & \textbf{Timeline UI:} Four temporal resolutions (day, week, month, year), LLM-generated per-week, per-month, per condition summaries \textcolor{phendopink}{(Timeline~1--4; Reflective Sensemaking~1--2)}.                                                                                                                                       \\
    \hline
    Patients spend hours preparing for visits \& need to construct narrative \& evidence and control what is disclosed (Interviews; \cite{Pichon2021Divided}).           & \textbf{DG4.} Produce a visit-preparation artifact the patient controls.                    & \textbf{Visual Summarization UI:} Anterior/posterior body map with tappable/clickable severity-colored zones, four lookback time ranges, narrative summary per range, patient controlled disclosure \textcolor{phendopink}{(Visit Prep~1--6)}.                                                                                                   \\
    \hline
    Patients fear data exposure and have privacy concerns.                                                                                                               & \textbf{Cross-cutting.} Preserve privacy as precondition for trust.                         & Exclusively on-device storage, no telemetry, audio discarded after transcription, deidentified transcripts sent upstream for summaries only \textcolor{phendopink}{(Identity \& Privacy~1--4)}.                                                                                                                                                  \\
    \hline
    Patients risk harm if tracking becomes another source of pressure or failure\cite{McKillop2019Thesis}.                                                               & \textbf{Cross-cutting.} Do No Harm in design of intervention.                               & No gamification or pressure to log consistently \textcolor{phendopink}{(Capture~8)}, no language implying clinical certainty or medical advice \textcolor{phendopink}{(Reflective Sensemaking~3)}.                                                                                                                                               \\
    \hline
  \end{tabular}
  \caption{Mapping from User Needs $\rightarrow$ Design Goals $\rightarrow$ Features. \textcolor{phendopink}{Feature numbers} correspond to the Project Requirements in Appendix~\ref{sec:requirements}.}
  \label{tab:needs-goals-features}
\end{table}

\subsection{Prototype}

\subsubsection{Architecture and Implementation}

Figure~\ref{fig:architecture} shows the overall architecture of our intended solution. An interactive prototype of EndoChron is deployed at \href{https://docker.nikhil.io}{\texttt{docker.nikhil.io}} with the source available on \href{https://github.com/afreeorange/endochron}{GitHub}. The prototype implements all four design goals across three views: \textit{Record}, \textit{Reflect}, and \textit{Prepare}. Bruce Tognazzini's \textit{
  First Principles of Interaction Design}\cite{tognazzini2014interaction} were consulted prior to initial design and functioned as a cohesive set of principles throughout the iterative design process.

\begin{figure}[H]
  \centering
  \makebox[\textwidth][c]{%
    \includegraphics[width=1.2\textwidth]{./media/Architecture.jpg}
  }
  \caption{\small Proposed EndoChron Architecture. Users speak freely about symptoms (1), then curate, synthesize, and reflect on AI-generated timelines and body-map summaries (2) before optionally sharing with clinicians or caregivers (3). On-device transcription and vocabulary mapping preserve privacy. De-identified cloud-based summarization and health-rating models support longitudinal synthesis.}
  \label{fig:architecture}
\end{figure}

\subsubsection{Synthetic Data}

At this time, Synthetic Data exclusively drives the prototype (i.e., no audio is actually captured.) The data is meant to simulate a fictitious patient's account of living with endometriosis over a 15-month period. It was generated using Claude Opus 4.7 using the prompt shown in Appendix \ref{sec:prompt}. It may also be examined \href{https://github.com/afreeorange/endochron/blob/main/src/data/syntheticData.json}{on Github}. We simulate $\approx$440 days of app use with a 65\% engagement rate.

\subsubsection{User-Centered Design}

Using an initial prototype, we engaged in iterative, user-centered design, with multiple `live-coding' feedback sessions using the same pool of participants involved in the requirements gathering phase. The current version of EndoChron incorporates almost all UI and UX critique and feedback received during these sessions. Appendix \ref{fig:iphone-screens} shows some screenshots of the final interactive prototype.

% =============================================================

\section{Evaluation}

Our evaluation plan proceeds in three phases. Phase~I (formative, completed) used iterative usability sessions with an interactive prototype to refine the UI and UX. Phase~II (proposed) will be a Randomized Controlled Trial (RCT) designed to test the clinical \textit{efficacy} of a \textit{fully-functional version} of EndoChron. We plan on Phase~III being a future \textit{effectiveness} trial.

\subsection{Phase I: Formative Design and Usability Evaluation (Completed)}

All five participants completed individual semi-structured usability sessions with the interactive prototype. Sessions used a think-aloud protocol\cite{Ericsson1980Verbal} and lasted $\approx$30 minutes on average. The Record, Reflect, and Prepare views were visited in order, with the participant asked to imagine themselves as the synthetic patient preparing for a clinical visit. All flagged bugs and most UX concerns or suggestions were noted and integrated immediately into the prototype via `live-coding' sessions where each participant witnessed their feedback reflected in the app in near real-time.

\subsection{Phase II: Clinical Efficacy RCT (Proposed)}
\label{sec:efficacy-rct}

In this phase, we will build an initial, fully-functional version of EndoChron per the design requirements identified in Phase I and described in Table \ref{tab:needs-goals-features}. Our research question would be: \textit{Among Phendo users with Endometriosis, does adding EndoChron for eight weeks preceding an OB/GYN visit improve patient-reported preparedness and self-efficacy, compared to Phendo alone?}

We will note that \textbf{we are not contrasting the efficacies of Phendo and EndoChron in this design}. For one, we would be comparing two different self-tracking modalities and not isolating EndoChron's contribution (if any.) We also wish to avoid ethical and recruitment issues: Phendo users are recruited into Phendo for a research study. Telling them to cease using Phendo for 8–12 weeks and use EndoChron instead asks them to drop out of a study they are enrolled in. We believe that our design is a protocol \textit{amendment} since both arms are still comprised of Phendo users. Comparing Phendo to EndoChron might necessitate a separate protocol.

Patient-reported visit preparedness and self-efficacy would measured by the Perceived Efficacy in Patient--Physician Interactions Scale\cite{Maly1998PEPPI,tenKlooster2012PEPPI} (PEPPI-5), a validated and brief instrument with strong psychometric properties in chronic disease populations\cite{tenKlooster2012PEPPI}. The instrument is conceptually aligned with EndoChron’s grounding in Distributed Cognition because EndoChron is explicitly designed to reduce cognitive burden during clinical encounters by transforming fragmented lived experiences into collaboratively referenced, persistent external representations. PEPPI-5 will be administered within 24~hours \textit{before} the index visit and is the primary outcome of interest.

Our Null hypothesis is that there will be no difference in patient-reported preparedness and self-efficacy between patients who use Phendo with EndoChron and patients who only use Phendo.

Our Alternative hypothesis is that \textit{patients who use EndoChron will arrive at their clinical visit with \circled{1} a more complete representation of their longitudinal experience, \circled{2} greater perceived self-efficacy in patient--provider interaction, and \circled{3} higher satisfaction with communication post-visit, than patients using Phendo alone.}

Statistically, we hypothesize that \textit{the mean PEPPI-5 score at the index visit will be higher in those using both Phendo and EndoChron than in those using just Phendo, by Cohen's} $d \geq 0.5$.

\subsubsection{Design}

We propose a two-arm RCT with 1:1 allocation and patient-level randomization. The design is a \textit{superiority trial}\cite{kishore2020superiority} in that both arms continue usual Phendo use but the intervention arm additionally receives EndoChron. This isolates the marginal effect of EndoChron beyond the self-tracking baseline and rules out attributing any effect to attention or to general self-tracking. Randomization will be stratified on the level of Phendo use/engagement (Regulars and Usuals vs.\ Occasionals and Seldoms\cite{Urteaga2022Engagement})\footnote{We argue that this \textit{may} tenably be associated with the primary outcome and is ascertainable at baseline from the Phendo usage data.}.

\textit{Intervention Arm.} EndoChron is installed on the participant's device alongside Phendo. Participants are encouraged but not required to use the app. The eight-week period preceding the index visit is the active intervention window. In the 48~hours preceding the visit, the intervention is augmented with a structured nudge to complete the Prepare view, but use remains the participant's choice (consistent with DG2 and the `no coercion' requirement.)

\textit{Control Arm.} Continued use of Phendo as before, with no access to EndoChron. To address the Hawthorne Effect\cite{Hutchins1995Cognition}, control participants will complete equivalent baseline and outcome instruments.

\subsubsection{Participants, Sample Size, Power, and Analysis}

We plan on recruiting participants through the Phendo research panel and Citizen Endo community, under the existing IRB protocol \texttt{\#AAAQ9812}\cite{Urteaga2022Engagement}. We will \textbf{include} adults 18+ with confirmed endometriosis, $\geq$4 months of Phendo use, a scheduled OB/GYN visit within 8--12 weeks of enrollment, a modern smartphone, English proficiency for the speech interface and PEPPI-5. We will \textbf{exclude} those with active hospitalization for endometriosis or prior EndoChron usability participation. All participants and will receive the same compensation. EndoChron will be offered to controls after the index visit.

Studies have reported PEPPI-5 standard deviations ranging from 5 to 8 points, depending on the scoring scale and study population \cite{raymond2011selfefficacy,tenKlooster2012PEPPI}. Assuming a conservative standard deviation of $\sigma = 8.0$ on the standard scale, a clinically relevant raw difference of 4.0 points between groups corresponds to a medium effect size (Cohen's $d \gtrsim 0.5$). To detect this with a two-sided significance level of $\alpha = 0.05$ and a power of $1 - \beta = 0.80$, an independent samples $t$-test requires $n = 64$ evaluable participants per arm. Factoring in an anticipated 20\% attrition or non-participation rate ($\frac{n}{0.8}$), the initial enrollment target is $n = 80$ per arm, yielding a total required sample size of $N = 160$. This target is highly feasible within Phendo's existing research infrastructure ($N = 4{,}993$ as of 2022 \cite{Urteaga2022Engagement}).

The primary analysis will be a comparison of mean PEPPI-5 score between arms using a two-sample Welch's $t$-test\cite{Welch1947Generalization} which is robust to unequal variances. A secondary ANCOVA model will include the baseline PEPPI-5 score and engagement strata (Regulars and Usuals vs.\ Occasionals and Seldoms) as covariates to improve precision and assess potential differences in response across baseline engagement levels.

A secondary sensitivity analysis \textit{might} examine whether the effect varies by engagement level which would require a larger sample. At this time, we will deem this analysis tentative, exploratory, and hypothesis-generating.

\subsubsection{Outcomes}

Our \textit{Primary} outcome of interest is the PEPPI-5 score, administered within 24~hours before the index visit. Our \textit{Secondary} outcomes may be:

\begin{itemize}
  \item \textbf{Post-visit Communication Quality:} the CAHPS Clinician \& Group Survey communication composite\cite{Dyer2012CAHPS}, administered 72~hours after the index visit.
  \item \textbf{Patient Autonomy:} the Autonomy subscale of the Health Care Climate Questionnaire\cite{Williams1996Autonomy}, grounded in Self-Determination Theory (SDT). This outcome is motivated by the SDT framing in original Phendo motivations research\cite{McKillop2016Participatory}\footnote{``\textit{This 15-item scale assesses participants' perceptions of the degree of autonomy support (vs. controllingness) of the relevant health-care providers. It includes items such as “I feel that the staff has provided me with choices and options.” Answers are rated on a 5-point scale ranging from not true at all to very true. The HCCQ has an alpha of .95 based on a sample of 276 patients who visited a Rochester-area internal medicine office.}''}.
  \item \textbf{Provider-Rated Information Quality:} a brief 4-to-5 item questionnaire based on the Coiera Common-Ground Framework\cite{Coiera2000ConversationModel}, asking the provider (blinded to arm allocation) to rate the completeness, clarity, and clinical utility of the information the patient brought to the visit (may also capture  proportion of visit time spent on shifting-ground work.)
\end{itemize}

\newpage
With regard to \textit{Process \& Safety},

\begin{itemize}
  \item \textbf{EndoChron Engagement Metrics}: Number of recordings, total speech duration, number of corrections, body-map zones annotated, \textit{Prepare}-view completions. Pushed from the participant's device and with consent (no telemetry is a design goal.)
  \item \textbf{Adverse Event Reporting:} A dedicated phone line and email address will be provided for participants to report distress, intrusive worry, or other negative affect attributable to using the app. Reports will be triaged by the research team and referred to clinical support appropriately.
\end{itemize}

\subsubsection{Threats to Validity and Mitigations}

Some threats to \textit{Internal Validity} may be:

\begin{itemize}
  \item \textit{Selection Bias.} Phendo users are not representative of all endometriosis patients. They are younger and more health-literate\cite{Urteaga2022Engagement}. This limits generalizability but does not threaten internal validity given their randomization in Phase II.
  \item \textit{Hawthorne Effect.} Both arms receive equivalent attention and instruments. Provider-rated outcomes will be collected with blinding to arm where feasible (note that the provider sees the patient's prepared artifact, so full blinding is not possible.)
  \item \textit{Assessment Bias.} As the developers of EndoChron, we will not be involved in outcome data collection, consistent with Friedman \& Wyatt's recommendation that developers should strive to not evaluate their own systems\cite{Friedman2006EvaluationMethods}.
\end{itemize}

We foresee no threats to \textit{External Validity} at this phase since the trial is confined to existing Phendo users, an inherently engaged population. Phase~III is designed to address this.


\subsection{Phase III: Effectiveness and Implementation (Proposed)}

Should Phase~II demonstrate efficacy, Phase~III will be a pragmatic \textit{effectiveness} study using the RE-AIM framework\cite{Glasgow1999REAIM} to evaluate Reach, Effectiveness, Adoption, Implementation, and Maintenance in a real-world deployment scenario beyond the Phendo research panel. We believe that it is the appropriate framework for Phase~III evaluation since it addresses external validity, organizational adoption, and sustainability over time, questions that we cannot answer with a strict efficacy-seeking/focused RCT.

% =============================================================

\newpage
\section{Human Subjects, IRB, and Data}
\label{sec:human-subjects}
\textbf{Whom We will Involve.} Phase~I (complete) involved five participants assigned female sex at birth, three with confirmed endometriosis. Sessions used semi-structured think-aloud usability methods with the interactive prototype and the synthetic patient dataset described in Appendix~D. Phase~II (proposed) will enroll $N=160$ adult Phendo users (80 per arm) with confirmed endometriosis and a scheduled gynecology visit, recruited through the existing Phendo and Citizen Endo recruitment infrastructure. Phase~III is outlined but not the current commitment.

\textbf{Risks and Benefits.} Reflecting on chronic illness can bring up difficult emotions: pain, frustration, grief, and memories of dismissal by providers. Sharing health information also carries privacy risks, especially for a condition that is stigmatized and historically dismissed. Anticipated benefits include the creation of a space in which the patient's account is treated as authoritative; support for the cognitive and emotional preparation that patients already undertake before clinical visits; and, if Phase~II is successful, evidence that conversational-speech-based representations are a feasible alternative to structured-form self-tracking for an enigmatic disease.

\textbf{Participant Protections.} Participants will provide informed consent and may withdraw at any time without penalty. EndoChron stores all health data on-device; patients may delete it at any time by clearing local storage or uninstalling the app. Aggregate engagement metrics are shared only upon consent. Audio is discarded after transcription. Data submitted to upstream LLMs for summarization will be free of personally identifying information: requests will \textit{only entail transcripts} and not include name, age, gender, location, device signatures, or any other personally identifiable information. For Phase~II, an adverse event surveillance item will be administered at the post-visit timepoint asking whether the app caused distress, intrusive worry, or other negative affect. Phendo engagement data accessed for stratification will be governed by the existing IRB-approved data-use agreement. Phendo stratification data are accessed under the existing data-use agreement and de-identified.

\textbf{IRB.} All procedures will be reviewed by the Columbia IRB under an amendment to the existing Phendo protocol (\href{https://pmc.ncbi.nlm.nih.gov/articles/PMC7314826/}{\texttt{AAAQ9812}}), consistent with prior Phendo research.

\textbf{Inclusion.} Endometriosis primarily affects people with female reproductive anatomy. We will recruit women and welcome non-binary, transgender, and intersex people with the condition. Phase~II will exclude participants under 18, consistent with the existing Phendo protocol's pediatric provisions. Race, ethnicity, education, and socioeconomic status will be measured at baseline and reported transparently. The existing Phendo cohort is known to over-represent white, college-educated, English-speaking participants\cite{Urteaga2022Engagement}. Phase~III is planned in part to address that limitation.

\textbf{Broader Ethics.} Endometriosis patients are routinely dismissed by the healthcare system and historically under-studied. Following the citizen-science approach of prior Phendo work\cite{McKillop2016Participatory}, we will treat participants as collaborators rather than subjects, center their interpretations as the gold standard against which algorithmic representations are checked, and share findings back with the patient community. EndoChron is intentionally designed to complement clinical care and human relationships, not to replace them. No chatbot or simulated-empathy features are included. No attempt is made to override participants' narratives at any temporal resolution. We position the summarization AI as a reflective anchor and explicitly state this in the UI with explicit calls-to-action that all AI-generated content is editable by the user at any time.


\newpage


% =============================================================
\newpage
\renewcommand*{\bibfont}{\footnotesize}
\printbibliography

% =============================================================

\newpage
\section{Appendix}
\subsection{Project Requirements}
\label{sec:requirements}

\begingroup
\fontsize{10pt}{12pt}\selectfont
\ttfamily
\setlength{\parindent}{0pt}
\setlength{\parskip}{6pt}

The verbs ``MUST,'' ``MUST NOT,'' ``SHOULD,'' ``SHOULD NOT,'' and ``MAY'' are to be interpreted per \href{https://datatracker.ietf.org/doc/html/rfc2119}{RFC~2119}.

\textbf{Identity Management and Privacy}

\begin{enumerate}
  \item The system MUST present the option to sign in with Phendo credentials. It MUST also present the option to create a profile for patients who do not use Phendo.
  \item The system MUST inform the patient that all data is stored on-device via a dedicated privacy page.
  \item The system MUST NOT transmit, collect, or share any health data. No server-side storage, analytics, crash reports, or device identifiers SHALL be present.
  \item The system MUST allow the patient to delete all data by clearing local storage or uninstalling the app.
\end{enumerate}

\textbf{Capture}

\begin{enumerate}
  \item The system MUST allow the patient to record their experience through open-ended speech via a dedicated recording interface with a visual waveform. It MUST NOT impose structured input.
  \item The system MUST state that all audio is deleted after transcription.
  \item The system SHOULD use an on-device LLM for speech transcription. (The prototype demonstrates the transcription flow with synthetic data; live on-device transcription is deferred.)
  \item The system MUST extract clinically relevant factors from conversational input using a downstream LLM which MAY be off-device. Extracted factors MUST include:
        \begin{itemize}
          \item Pain (16 locations, 3 severity levels)
          \item Mood (10 positive, 20 negative sentiments)
          \item Period (flow intensity, clots, spotting, breakthrough bleeding)
          \item GI/Urinary symptoms (15 symptoms, 3 severity levels)
          \item Functional impact (``Hard to Do'': 20 activities)
          \item Other symptoms (21 items, 3 severity levels)
          \item Medications (free-text list)
          \item Overall day rating (Good / Manageable / Bad)
        \end{itemize}
  \item Extraction categories MUST align with the Phendo citizen-science taxonomy.
  \item The system MUST surface all extractions transparently as visual badges/pills organized by category so the patient can see what was inferred.
  \item The system MUST allow the patient to override the overall day rating and edit the transcript text. It SHOULD allow the patient to correct or reject individual extracted factors.
  \item The system MUST NOT gamify tracking or create pressure to log consistently.
\end{enumerate}

\textbf{Reflective Sensemaking}

\begin{enumerate}
  \item The system MUST generate per-week and per-month narrative summaries covering pain, mood, period, GI, functional impact, medications, and overall status.
  \item All AI-generated summaries MUST be presented as reflective anchors---observational and patient-centered in tone (e.g., ``Here's what I see in your data'').
  \item The system MUST NOT generate language that implies clinical certainty or medical advice.
  \item Weekly summaries MUST be viewable by month, with collapsible long/short display modes.
  \item Monthly summaries MUST aggregate symptom frequencies and top factors.
  \item The patient SHOULD be able to edit summaries inline.
\end{enumerate}

\textbf{Timeline}

\begin{enumerate}
  \item The system MUST support viewing the illness trajectory at four temporal resolutions: day, week, month, and year.
  \item Day view MUST display a horizontally scrollable date strip organized by month, with each day showing the overall rating and full detail on selection.
  \item Year view MUST render a calendar heatmap colored by category severity.
  \item Year view MUST allow the patient to switch the displayed category among Overall, Pain, Mood, GI, and Period. The selected category SHOULD persist across sessions.
  \item The system SHOULD support arbitrary, user-defined timespans. (Not yet delivered.)
  \item The system SHOULD allow the patient to annotate any point on the timeline with free text and tags. (Not yet delivered.)
\end{enumerate}

\textbf{Visit Preparation}

\begin{enumerate}
  \item The system MUST provide a clinical visit preparation view with interactive anterior and posterior body map visualizations.
  \item Body map zones MUST be clickable and MUST display a summary of mapped symptoms, pain locations, and their frequencies for the selected time range.
  \item Zones MUST be color-coded by maximum severity within the selected time range.
  \item The system MUST support at least four preparation time ranges: last week, two weeks, last month, and six months.
  \item Each time range MUST present a narrative summary covering logging frequency, pain days/locations, period data, GI patterns, functional impact, medications, and mood themes.
  \item The system SHOULD allow the patient to control which data is disclosed to a clinician.
\end{enumerate}

\endgroup

\subsection{EndoChron Prototype Screenshots}
\label{sec:screenshots}

\begin{sidewaysfigure}[!htbp]
  \centering
  \captionsetup[subfigure]{font=scriptsize}

  \begin{subfigure}{0.23\linewidth}
    \centering
    \includegraphics[height=0.35\textheight]{./media/Screenshots/17_08.png}
    \caption{User authentication using either Phendo credentials or a new profile}
  \end{subfigure}
  \hfill
  \begin{subfigure}{0.23\linewidth}
    \centering
    \includegraphics[height=0.35\textheight]{./media/Screenshots/17_15.png}
    \caption{The Recording UI}
  \end{subfigure}
  \hfill
  \begin{subfigure}{0.23\linewidth}
    \centering
    \includegraphics[height=0.35\textheight]{./media/Screenshots/17_28.png}
    \caption{The Recording UI plus live transcription}
  \end{subfigure}
  \hfill
  \begin{subfigure}{0.23\linewidth}
    \centering
    \includegraphics[height=0.35\textheight]{./media/Screenshots/17_36.png}
    \caption{Reflect UI: Daily View. All entries are editable.}
  \end{subfigure}

  \vspace{0.8em}

  \begin{subfigure}{0.23\linewidth}
    \centering
    \includegraphics[height=0.35\textheight]{./media/Screenshots/17_45.png}
    \caption{Reflect UI: Weekly View.}
  \end{subfigure}
  \hfill
  \begin{subfigure}{0.23\linewidth}
    \centering
    \includegraphics[height=0.35\textheight]{./media/Screenshots/17_48.png}
    \caption{Reflect UI: Monthly View.}
  \end{subfigure}
  \hfill
  \begin{subfigure}{0.23\linewidth}
    \centering
    \includegraphics[height=0.35\textheight]{./media/Screenshots/17_51.png}
    \caption{Reflect UI: Yearly View.}
  \end{subfigure}
  \hfill
  \begin{subfigure}{0.23\linewidth}
    \centering
    \includegraphics[height=0.35\textheight]{./media/Screenshots/18_07.png}
    \caption{Prepare UI, with editable conditions annotated on a Body Map}
  \end{subfigure}

  \caption{Representative EndoChron prototype screens across the \textit{Record}, \textit{Reflect}, and \textit{Prepare} views.}
  \label{fig:iphone-screens}
\end{sidewaysfigure}

\newpage
\input{sections/Appendix/Prompt.tex}
\newpage
\subsection{Phendo Taxonomy}
\label{sec:taxonomy}

\vspace{0.75em}

\begingroup
\fontsize{8pt}{1pt}\selectfont
\ttfamily
\setlength{\parindent}{0pt}
\setlength{\parskip}{6pt}

\begin{multicols}{2}
  \textbf{\uppercase{Signs \& Symptoms}}

  \vspace{0.125em}
  \noindent\rule{\columnwidth}{0.25pt}
  \vspace{0.0625em}

  \textbf{Daily Variables [day]}
  \begin{itemize}
    \item How was your day?
  \end{itemize}

  \textbf{Period \& Bleeding [day/moment]}
  \begin{itemize}
    \item Period Y/N (heavy, med, light) [day]
    \item Clots [moment]
    \item Spotting [moment]
    \item Breakthrough Bleeding [moment]
  \end{itemize}

  \textbf{Pain [moment] (mild/mod/severe)}
  \begin{itemize}
    \item ALL
    \item Pelvis (W/L/R)
    \item Ovary (L/R)
    \item Uterus
    \item Vagina Entrance
    \item Deep Vagina
    \item Cervix
    \item Rectum
    \item Lower Back (W/L/R)
    \item Outer Hip (L/R)
    \item Abdomen (W/L/R)
    \item Head
    \item Neck
    \item Inner Thighs
    \item Joints
    \item Bones
    \item Shoulder (L/R)
    \item Ribs (L/R)
    \item Breast (L/R)
    \item Arm (L/R)
    \item Upper Chest
    \item Lower Chest
    \item Upper Abdomen
    \item Diaphragm
    \item Intestines
    \item Leg (W/L/R)
  \end{itemize}

  \textbf{Describe the Pain}
  \begin{itemize}
    \item Aching
    \item Burning
    \item Cramping
    \item Deep
    \item Dull
    \item Pressure
    \item Pulling
    \item Radiating
    \item Sharp
    \item Shooting
    \item Stabbing
    \item Throbbing
    \item Twisting
    \item Pulsating
    \item Nauseating
  \end{itemize}

  \textbf{GI/Urinary [moment] (mild/mod/severe)}
  \begin{itemize}
    \item ALL
    \item Nausea
    \item Endo Belly
    \item Stomach Upset
    \item Vomiting
    \item Diarrhea
    \item Constipation
    \item Uncomfortably Full
    \item Heartburn
    \item Gas
    \item Mouth Sores
    \item Can't Urinate
    \item Painful Urination
    \item Painful Bowel Movement
    \item Frequent Urination
    \item Blood in Stool
  \end{itemize}

  \textbf{Something Else [moment] (mild/mod/severe)}
  \begin{itemize}
    \item ALL
    \item Fatigue
    \item Numbness
    \item Headache
    \item Asthma
    \item Chest Pressure
    \item Swelling
    \item Rash
    \item Ringing in Ears
    \item Eczema
    \item Hives
    \item Allergies
    \item Itchy
    \item Hot Flash
    \item Sweaty
    \item Touch Sensitivity
    \item Noise Sensitivity
    \item Mentally Foggy
    \item Sinus Congestion
    \item Fever
    \item Dizziness
    \item Blurry Vision
  \end{itemize}

  \textbf{Mood [moment]}\\
  \textit{Positive}
  \begin{itemize}
    \item ALL positive
    \item Affectionate
    \item Calm
    \item Excited
    \item Enthusiastic
    \item Happy
    \item Motivated
    \item Optimistic
    \item Productive
    \item Relaxed
    \item Social
  \end{itemize}
  \textit{Negative}
  \begin{itemize}
    \item ALL negative
    \item Angry
    \item Antisocial
    \item Anxious
    \item Belligerent
    \item Contemptuous
    \item Erratic
    \item Defensive
    \item Disgusted
    \item Frustrated
    \item Guilty
    \item Indifferent
    \item Irritable
    \item Isolated
    \item Lonely
    \item Mentally Foggy
    \item Overwhelmed
    \item Sad
    \item Scared
    \item Whiny
    \item Worried
  \end{itemize}

  \textbf{Sex [day]} (penetration: vaginal, anal, none)
  \begin{itemize}
    \item ALL
    \item Painful During
    \item Painful After
    \item Felt Good
    \item Avoided
    \item Bleeding from Sex
  \end{itemize}

  \textbf{Activities Hard to Do [day]}
  \begin{itemize}
    \item ALL
    \item Sleep
    \item Get out of Bed
    \item Use Toilet
    \item Shower
    \item Get Dressed
    \item Prepare Food
    \item Eat
    \item Sit Down
    \item Work
    \item Stand
    \item Stretch
    \item Socialize
    \item Shop
    \item Have Sex
    \item Lie Down
    \item Run
    \item Walk
    \item Jump
    \item Climb Stairs
    \item Kneel
  \end{itemize}

  \vspace{1em}

  \textbf{\uppercase{Self Management}}

  \vspace{0.125em}
  \noindent\rule{\columnwidth}{0.25pt}
  \vspace{0.0625em}

  \textbf{Strategies [day]} (didn't help / no effect / helped)
  \begin{itemize}
    \item ALL
    \item Breathing Exercises
    \item Heat Pack
    \item Ice Pack
    \item Stretching
    \item Massage
    \item TENS
    \item Pelvic Floor Therapy
    \item Physical Therapy
    \item Talk Therapy
    \item Acupuncture
    \item Medical Marijuana
    \item Alcohol
    \item CBD Oil
    \item Rest
  \end{itemize}

  \textbf{Hormones [day]}
  \begin{itemize}
    \item ALL
    \item \ldots
  \end{itemize}

  \textbf{Medications [moment]}
  \begin{itemize}
    \item ALL
    \item \ldots
  \end{itemize}

  \textbf{Supplements [day]}
  \begin{itemize}
    \item ALL
    \item \ldots
  \end{itemize}

  \textbf{Foods [day]}\\
  \textit{hurt}
  \begin{itemize}
    \item ALL
    \item \ldots
  \end{itemize}
  \textit{help}
  \begin{itemize}
    \item ALL
    \item \ldots
  \end{itemize}

  \textbf{Exercise [day]} (mild / mod / severe / no effect)\\
  \textit{hurt}
  \begin{itemize}
    \item ALL
    \item \ldots
  \end{itemize}
  \textit{help}
  \begin{itemize}
    \item ALL
    \item \ldots
  \end{itemize}

  \textbf{Time}\\
  \textit{what to show}
  \begin{itemize}
    \item Week
    \item Month
    \item Cycle
  \end{itemize}
  \textit{how many}
  \begin{itemize}
    \item 1
    \item 3
    \item ALL
  \end{itemize}

  \textbf{Compare To}
  \begin{itemize}
    \item All Phendo Users
    \item My geographical area
    \item My Endo
  \end{itemize}

  \textbf{Metadata}
  \begin{itemize}
    \item Days/Moments Tracked (use lvl)
    \item Future: geofence, weather info, steps/activity
  \end{itemize}

\end{multicols}


\end{document}
